%----------
%	EXPERIMENTACIÓN Y RESULTADOS (versión de trabajo)
%
%	Copia de experimentacion.tex sobre la que se aplican los comentarios. El original se conserva intacto hasta que esta versión
%	quede cerrada y se cambie el \input del fichero principal.
%----------

% =====================================================================
% Experimentación y resultados — Agentes especializadas y Agent Skills
% Datos: Experimentos/experiments/experimento_21_agentskills/results/
% Corpus de evaluación: 1.313 noticias anotadas en las cinco variables
% =====================================================================

La experimentación se organiza en torno a tres preguntas. La primera es si aporta
algo empaquetar la metodología experta como \textit{Agent Skills}, para lo cual se
enfrenta la arquitectura propuesta (\textbf{nivel B1}) a una condición de control
que recibe la misma metodología por la vía más simple, inyectada íntegramente en el
\textit{prompt} (\textbf{nivel B0}). Ambos niveles se definen en el
Capítulo~\ref{implementation} (Sección~\ref{subsec:diseno_agentico}). La segunda
pregunta atiende a los modelos, esto es, cuáles se aproximan más al criterio
experto, cuáles menos y a qué obedecen esas diferencias. La tercera atraviesa a las
dos anteriores y se refiere a las variables, ya que cada una plantea una tarea de
detección distinta y nada garantiza que una misma configuración resulte igual de
adecuada para todas.

Estas tres preguntas no admiten una respuesta única. Como se verá, el efecto del
empaquetado depende del modelo, el rendimiento de cada modelo depende de la
variable, y la configuración más conveniente puede diferir de una variable a otra.
El análisis se organiza en consecuencia, variable por variable, y solo al final se
valora hasta qué punto cabe extraer alguna conclusión de conjunto.

Como se detalla en la Sección~\ref{subsec:diseno_agentico}, el diseño mantiene fijo
todo salvo el empaquetado (mismo corpus, modelos, variables y formato de salida),
de modo que cualquier variación en el acuerdo es atribuible a cómo llega la
metodología al modelo y no a otro factor. Para homogeneizar además el comportamiento
entre empresas proveedoras, se desactivó el razonamiento explícito en los modelos que lo
admiten (Sección~\ref{sec:genai_models_iris}). La
interpretación de las métricas exige tener presente la prevalencia de cada variable
y la exactitud trivial asociada, recogidas en la
Tabla~\ref{tab:iris-prevalencia-imp} del Capítulo~\ref{sec:methodology}.

\subsection{Panorámica: el efecto del empaquetado en cada modelo}
\label{subsec:iris_core}

Antes de descender al detalle por variable conviene una panorámica que sitúe a los
cuatro modelos y ambos niveles. La Tabla~\ref{tab:iris-b0b1} los compara en los
cuatro modelos evaluados, los tres accesibles por API de pago y el modelo local
\texttt{gemma4:e4b}. Sus valores son promedios sobre las cinco variables, de modo
que resumen la posición de cada modelo pero ocultan diferencias considerables entre
variables, según se comprueba en la Sección~\ref{subsec:iris_porvariable}. Conviene
leerlos, por tanto, como una primera orientación y no como el resultado del
capítulo.

\begin{table}[H]
    \centering
    \begingroup
    \scriptsize
    \setlength{\tabcolsep}{4pt}
    \newcommand{\dr}[1]{\textcolor{gray!75}{$#1$}} % fila delta: en gris
    \renewcommand{\texttt}[1]{{\ttfamily #1}}
    \ttabbox[\FBwidth]{
        \caption{Comparación B0 (\textit{baseline}) frente a B1 (Agent Skills). Medias sobre las cinco variables, $N=1.313$. Las filas $\Delta$ (en gris) recogen la variación B1$-$B0.}
        \label{tab:iris-b0b1}
    }{
        \begin{tabularx}{\textwidth}{
            >{\hsize=1.7\hsize\raggedright\arraybackslash}X
            >{\hsize=0.6\hsize\centering\arraybackslash}X
            >{\hsize=0.925\hsize\centering\arraybackslash}X
            >{\hsize=0.925\hsize\columncolor{gray!12}\centering\arraybackslash}X
            >{\hsize=0.925\hsize\centering\arraybackslash}X
            >{\hsize=0.925\hsize\centering\arraybackslash}X
        }
            \toprule
            \textbf{Modelo} & \textbf{Nivel} & \textbf{Exactitud} & \boldmath$\kappa$ & \textbf{F1 macro} & \textbf{Coste (USD)} \\
            \midrule
            \multirow{3}{*}{\texttt{gemini-3.1-flash-lite}}
              & B0 & 0,600 & 0,026 & 0,388 & 5,64 \\
              & B1 & 0,661 & \textbf{0,066} & 0,443 & 13,16 \\
              & \dr{\Delta} & \dr{+0,061} & \dr{+0,040} & \dr{+0,055} & \dr{2{,}3\times} \\
            \midrule
            \multirow{3}{*}{\texttt{gpt-4o-mini}}
              & B0 & 0,564 & 0,033 & 0,410 & 2,52 \\
              & B1 & 0,612 & 0,014 & 0,438 & 6,89 \\
              & \dr{\Delta} & \dr{+0,048} & \dr{-0,019} & \dr{+0,029} & \dr{2{,}7\times} \\
            \midrule
            \multirow{3}{*}{\texttt{gpt-5.4-nano}}
              & B0 & 0,571 & 0,052 & 0,430 & 5,67 \\
              & B1 & 0,569 & 0,004 & 0,364 & 10,86 \\
              & \dr{\Delta} & \dr{-0,003} & \dr{-0,048} & \dr{-0,066} & \dr{1{,}9\times} \\
            \midrule
            \multirow{2}{*}{\textbf{Media (API)}}
              & \textbf{B0} & 0,578 & \textbf{0,037} & 0,409 & \textemdash \\
              & \textbf{B1} & \textbf{0,614} & 0,028 & 0,415 & \textemdash \\
            \midrule
            \multirow{3}{*}{\texttt{gemma4:e4b}\,\textnormal{\footnotesize(local)}}
              & B0 & 0,597 & \textbf{0,061} & 0,421 & \textemdash \\
              & B1 & 0,605 & 0,043 & 0,412 & \textemdash \\
              & \dr{\Delta} & \dr{+0,008} & \dr{-0,018} & \dr{-0,009} & \textemdash \\
            \bottomrule
        \end{tabularx}
    }
    \endgroup
\end{table}

A primera vista, la arquitectura parece funcionar, ya que la exactitud sube con B1
en la mayoría de los modelos (Figura~\ref{fig:iris-b0b1-metricas}) y la exactitud
media de las API pasa de 0,578 a 0,614. Sin embargo, exactitud y $\kappa$ \emph{no
cuentan la misma historia}, puesto que el $\kappa$ medio, lejos de acompañar, baja
de 0,037 a 0,028, y todo ello a un coste de inferencia que se multiplica por un
factor de entre 1,9 y 2,7 según el modelo (Tabla~\ref{tab:iris-b0b1}). Esa
discordancia obliga a preguntarse qué está midiendo cada una. Con clases
fuertemente desequilibradas, la exactitud recoge en buena parte el efecto de la
distribución, de modo que un cambio en ella puede reflejar que el modelo marca con
otra frecuencia y no que discrimine mejor. El coeficiente $\kappa$ descuenta ese
componente, si bien tiene sus propias limitaciones
(Sección~\ref{sec:metricas_eval}). Ninguna de las dos basta por separado, y la
Sección~\ref{subsec:iris_divergencia} examina a qué responde esta divergencia.

\begin{figure}[H]
    \centering
    \includegraphics[width=\textwidth]{fig_b0b1_barras_blanco.pdf}
    \caption{Niveles de exactitud, F1 macro y $\kappa$ por modelo, en B0 y B1
    ($N=1.313$). Cada métrica se representa en su propia escala: el $\kappa$ vive
    en un orden de magnitud mucho menor que la exactitud.}
    \label{fig:iris-b0b1-metricas}
\end{figure}


\subsection{Divergencia entre exactitud y \texorpdfstring{$\kappa$}{kappa}}
\label{subsec:iris_divergencia}

Conviene detenerse en esa discordancia, porque condiciona cómo debe leerse toda la
tabla: \textbf{la mejora en exactitud no se traduce en más acuerdo}. Considerando
las medias, el nivel B1 sube la exactitud en 0,036 pero el $\kappa$ baja en $-0,009$.
Exactitud y $\kappa$ pueden, por tanto, moverse de forma independiente e incluso en
sentidos opuestos. Esto es un efecto conocido de $\kappa$ en
presencia de clases desequilibradas~\cite{feinstein1990}, pero aquí resulta determinante: pasarlo por alto
invertiría la conclusión del capítulo. La Figura~\ref{fig:iris-divergencia} lo
visualiza como trayectorias B0$\to$B1: en el panel de exactitud casi todas las líneas
ascienden, mientras que en el de $\kappa$ se mantienen o descienden, salvo en
\texttt{gemini}, cuya línea sube en ambas.

\begin{figure}[H]
    \centering
    \includegraphics[width=\textwidth]{fig_b0b1_slope_blanco.pdf}
    \caption{Evolución de cada modelo de B0 a B1 en las tres métricas ($N=1.313$).
    La exactitud sube en tres de los cuatro modelos, pero el $\kappa$ solo mejora en
    \texttt{gemini-3.1-flash-lite}. En el resto se mantiene o desciende. La mejora en
    exactitud sin mejora en $\kappa$ es el desacoplamiento entre exactitud y acuerdo.}
    \label{fig:iris-divergencia}
\end{figure}

El motivo de esa divergencia es conocido, y no exige que el modelo conozca la
distribución de clases: en régimen \textit{zero-shot} no influye. Al volverse menos
conservador (por efecto del empaquetado, no por una estrategia informada), el modelo
responde \textit{Sí} con más frecuencia. En las variables donde la clase positiva
resulta ser la mayoritaria, ese desplazamiento eleva la exactitud por simple
coincidencia con la mayoría, no por discriminar mejor. Se trata, por tanto, de un
artefacto estadístico de un sesgo ciego de respuesta, no de un mejor criterio. El
coeficiente $\kappa$, que descuenta el acuerdo esperable por azar, no premia ese
desplazamiento. Atender solo a la exactitud llevaría, por tanto, a dar por buena
una mejora que en parte responde a un cambio en la frecuencia de respuesta.

Conviene, con todo, no sustituir una lectura simplificada por otra, porque
$\kappa$ arrastra sus propias dificultades y este caso las reúne casi todas. La
primera es que el coeficiente se concibió para medir el acuerdo entre dos personas
codificadoras en pie de igualdad~\cite{cohen1960}, mientras que aquí una de las
partes es un patrón de referencia y la otra una estimación, situación para la que
no fue pensado. La segunda es que el acuerdo esperado por azar depende de las
frecuencias marginales de ambas partes, de manera que dos modelos con la misma
exactitud pueden obtener valores de $\kappa$ muy distintos solo por marcar con
frecuencias distintas. La tercera es la ya mencionada paradoja de
$\kappa$~\cite{feinstein1990}, que con prevalencias extremas hunde el coeficiente
aunque el acuerdo observado sea alto. De ello se sigue una consecuencia que
condiciona la estructura de este capítulo, y es que los valores de $\kappa$ de
variables con prevalencias tan dispares no son comparables entre sí ni admiten
promediarse con sentido, razón por la cual el análisis se realiza variable a
variable (Sección~\ref{subsec:iris_porvariable}) y las medias que aparecen en esta
sección deben tomarse como una orientación gruesa.

Lo que cabe sostener, en suma, es que exactitud, $\kappa$ y prevalencia responden a
preguntas distintas y que interpretarlas por separado conduce a conclusiones
opuestas, de modo que en tareas de clasificación desequilibradas procede informar
de las tres de forma conjunta.

Atendiendo al $\kappa$, que es la métrica que menos se deja arrastrar por la
frecuencia de marcado aunque no esté libre de problemas, el efecto de las
\textit{Agent Skills} \textbf{depende del modelo}. Solo \texttt{gemini-3.1-flash-lite} mejora con la
arquitectura: su $\kappa$ pasa de 0,026 a 0,066, y lo hace acompañado de subidas
simultáneas en exactitud y F1. En los dos modelos de OpenAI las \textit{skills} no ayudan e incluso
perjudican ($\Delta\kappa = -0,019$ en \texttt{gpt-4o-mini} y $-0,048$ en
\texttt{gpt-5.4-nano}), y el modelo local se comporta igual que ellos ($\kappa$ de
0,061 a 0,043). De los cuatro modelos, solo uno mejora con las \textit{skills}. En
los otros tres el $\kappa$ se mantiene o desciende, y la media no aumenta. Que el
efecto dependa del modelo y no eleve el acuerdo en agregado admite varias lecturas, desde el techo de la propia tarea hasta la heterogeneidad de la anotación, que se
retoman en la discusión (Sección~\ref{sec:discussion}) una vez presentada toda la
evidencia.

\subsection{Descomposición de la arquitectura por componentes (ablación)}
\label{subsec:iris_descomposicion}

El nivel B1 combina tres componentes, descritos en el Capítulo~\ref{implementation}:
la carga de \textit{skills} bajo demanda con verificación de evidencias
(Sección~\ref{subsec:diseno_agentico}), un catálogo de resúmenes de las guías
expertas descubrible por la agente (Sección~\ref{prompts}), y una herramienta de
recuperación en vivo sobre esas guías (RAG, implementada como \texttt{CONSULTAR\_GUIA},
Sección~\ref{rag}). Para averiguar cuál de ellos explica el efecto observado se
ejecutaron cuatro configuraciones por modelo, activando y desactivando cada
componente de forma independiente, sobre los cuatro modelos y las 1.313 noticias.

Las Tablas~\ref{tab:abl-v25} a~\ref{tab:abl-v35} recogen el resultado
\textbf{variable a variable}, con las cinco métricas de la
Sección~\ref{sec:metricas_eval} y el coeficiente $\kappa$. Se presentan así, y no
en agregado, porque las prevalencias van del 6,2\,\% al 81,0\,\% y los valores no
resultan comparables entre variables. En cada modelo se destaca su mejor valor por
métrica, ya que lo que aquí se compara son componentes dentro de un mismo modelo y
no modelos entre sí. Las escasas respuestas que el modelo no llegó a emitir en un
formato válido se contabilizan como negativas, de modo que todas las
configuraciones se evalúan sobre las mismas 1.313 piezas.

% Generado por gen_tablas_ablacion.py. No editar a mano.

\begin{table}[H]
    \centering
    \begingroup
    \small
    \setlength{\tabcolsep}{4pt}
    \ttabbox[\FBwidth]{
        \caption{Ablación por componentes en lenguaje sexista (V25), nivel B1 ($N=1.313$, prevalencia 0,810). Los modelos se nombran de forma abreviada, gemini por gemini-3.1-flash-lite y gemma por gemma4:e4b.}
        \label{tab:abl-v25}
    }{
        \begin{tabular*}{\textwidth}{@{\extracolsep{\fill}}l l c c c c c c@{}}
            \toprule
            \textbf{Modelo} & \textbf{Configuración} & \textbf{Exac.} & \textbf{Prec.} & \textbf{Recall} & \textbf{F1 (\textit{Sí})} & \textbf{F1 macro} & \textbf{$\kappa$} \\
            \midrule
            gemini & Solo \textit{skills} & 0,200 & 0,800 & 0,015 & 0,030 & 0,174 & 0,000 \\
             & \textit{Skills} + resúmenes & 0,215 & 0,971 & 0,031 & 0,060 & 0,193 & 0,011 \\
             & \textit{Skills} + RAG & 0,222 & 0,977 & 0,040 & 0,078 & 0,203 & 0,014 \\
             & Completo & 0,226 & 0,961 & 0,046 & 0,088 & 0,208 & 0,015 \\
            \midrule
            gpt-4o-mini & Solo \textit{skills} & 0,493 & 0,808 & 0,490 & 0,610 & 0,442 & $-$0,004 \\
             & \textit{Skills} + resúmenes & 0,458 & 0,806 & 0,437 & 0,566 & 0,423 & $-$0,007 \\
             & \textit{Skills} + RAG & 0,478 & 0,787 & 0,487 & 0,602 & 0,423 & $-$0,045 \\
             & Completo & 0,457 & 0,789 & 0,450 & 0,573 & 0,414 & $-$0,037 \\
            \midrule
            gpt-5.4-nano & Solo \textit{skills} & 0,192 & 1,000 & 0,002 & 0,004 & 0,162 & 0,001 \\
             & \textit{Skills} + resúmenes & 0,190 & 0,000 & 0,000 & 0,000 & 0,160 & 0,000 \\
             & \textit{Skills} + RAG & 0,190 & 0,000 & 0,000 & 0,000 & 0,160 & 0,000 \\
             & Completo & 0,191 & 1,000 & 0,001 & 0,002 & 0,161 & 0,000 \\
            \midrule
            gemma & Solo \textit{skills} & 0,198 & 0,857 & 0,011 & 0,022 & 0,171 & 0,001 \\
             & \textit{Skills} + resúmenes & 0,197 & 0,800 & 0,011 & 0,022 & 0,171 & 0,000 \\
             & \textit{Skills} + RAG & 0,194 & 0,727 & 0,008 & 0,015 & 0,167 & $-$0,002 \\
             & Completo & 0,201 & 0,938 & 0,014 & 0,028 & 0,175 & 0,004 \\
            \bottomrule
        \end{tabular*}
    }
    \endgroup
\end{table}

\begin{table}[H]
    \centering
    \begingroup
    \small
    \setlength{\tabcolsep}{4pt}
    \ttabbox[\FBwidth]{
        \caption{Ablación por componentes en masculino genérico (V26), nivel B1 ($N=1.313$, prevalencia 0,810). Los modelos se nombran de forma abreviada, gemini por gemini-3.1-flash-lite y gemma por gemma4:e4b.}
        \label{tab:abl-v26}
    }{
        \begin{tabular*}{\textwidth}{@{\extracolsep{\fill}}l l c c c c c c@{}}
            \toprule
            \textbf{Modelo} & \textbf{Configuración} & \textbf{Exac.} & \textbf{Prec.} & \textbf{Recall} & \textbf{F1 (\textit{Sí})} & \textbf{F1 macro} & \textbf{$\kappa$} \\
            \midrule
            gemini & Solo \textit{skills} & 0,645 & 0,910 & 0,624 & 0,740 & 0,590 & 0,237 \\
             & \textit{Skills} + resúmenes & 0,694 & 0,898 & 0,702 & 0,788 & 0,619 & 0,266 \\
             & \textit{Skills} + RAG & 0,662 & 0,891 & 0,664 & 0,761 & 0,592 & 0,224 \\
             & Completo & 0,708 & 0,887 & 0,733 & 0,803 & 0,621 & 0,261 \\
            \midrule
            gpt-4o-mini & Solo \textit{skills} & 0,644 & 0,824 & 0,713 & 0,765 & 0,518 & 0,051 \\
             & \textit{Skills} + resúmenes & 0,362 & 0,836 & 0,264 & 0,401 & 0,359 & 0,020 \\
             & \textit{Skills} + RAG & 0,506 & 0,827 & 0,494 & 0,619 & 0,460 & 0,032 \\
             & Completo & 0,332 & 0,873 & 0,206 & 0,333 & 0,332 & 0,034 \\
            \midrule
            gpt-5.4-nano & Solo \textit{skills} & 0,263 & 0,869 & 0,106 & 0,189 & 0,257 & 0,016 \\
             & \textit{Skills} + resúmenes & 0,253 & 0,919 & 0,086 & 0,156 & 0,243 & 0,021 \\
             & \textit{Skills} + RAG & 0,264 & 0,830 & 0,115 & 0,201 & 0,259 & 0,006 \\
             & Completo & 0,262 & 0,874 & 0,104 & 0,186 & 0,256 & 0,016 \\
            \midrule
            gemma & Solo \textit{skills} & 0,489 & 0,915 & 0,407 & 0,563 & 0,474 & 0,129 \\
             & \textit{Skills} + resúmenes & 0,492 & 0,908 & 0,415 & 0,570 & 0,475 & 0,124 \\
             & \textit{Skills} + RAG & 0,476 & 0,887 & 0,405 & 0,556 & 0,458 & 0,098 \\
             & Completo & 0,482 & 0,897 & 0,408 & 0,561 & 0,465 & 0,109 \\
            \bottomrule
        \end{tabular*}
    }
    \endgroup
\end{table}

\begin{table}[H]
    \centering
    \begingroup
    \small
    \setlength{\tabcolsep}{4pt}
    \ttabbox[\FBwidth]{
        \caption{Ablación por componentes en sexismo en el discurso (V30), nivel B1 ($N=1.313$, prevalencia 0,428). Los modelos se nombran de forma abreviada, gemini por gemini-3.1-flash-lite y gemma por gemma4:e4b.}
        \label{tab:abl-v30}
    }{
        \begin{tabular*}{\textwidth}{@{\extracolsep{\fill}}l l c c c c c c@{}}
            \toprule
            \textbf{Modelo} & \textbf{Configuración} & \textbf{Exac.} & \textbf{Prec.} & \textbf{Recall} & \textbf{F1 (\textit{Sí})} & \textbf{F1 macro} & \textbf{$\kappa$} \\
            \midrule
            gemini & Solo \textit{skills} & 0,571 & 0,467 & 0,012 & 0,024 & 0,375 & 0,002 \\
             & \textit{Skills} + resúmenes & 0,572 & 0,500 & 0,020 & 0,038 & 0,381 & 0,006 \\
             & \textit{Skills} + RAG & 0,567 & 0,406 & 0,023 & 0,044 & 0,382 & $-$0,002 \\
             & Completo & 0,569 & 0,423 & 0,020 & 0,037 & 0,380 & 0,000 \\
            \midrule
            gpt-4o-mini & Solo \textit{skills} & 0,558 & 0,446 & 0,141 & 0,214 & 0,453 & 0,011 \\
             & \textit{Skills} + resúmenes & 0,550 & 0,423 & 0,142 & 0,213 & 0,449 & $-$0,003 \\
             & \textit{Skills} + RAG & 0,561 & 0,456 & 0,139 & 0,213 & 0,454 & 0,016 \\
             & Completo & 0,561 & 0,450 & 0,112 & 0,179 & 0,440 & 0,011 \\
            \midrule
            gpt-5.4-nano & Solo \textit{skills} & 0,552 & 0,424 & 0,130 & 0,199 & 0,444 & $-$0,002 \\
             & \textit{Skills} + resúmenes & 0,566 & 0,469 & 0,107 & 0,174 & 0,440 & 0,018 \\
             & \textit{Skills} + RAG & 0,568 & 0,488 & 0,187 & 0,270 & 0,482 & 0,044 \\
             & Completo & 0,558 & 0,435 & 0,114 & 0,181 & 0,439 & 0,004 \\
            \midrule
            gemma & Solo \textit{skills} & 0,563 & 0,370 & 0,030 & 0,056 & 0,386 & $-$0,009 \\
             & \textit{Skills} + resúmenes & 0,570 & 0,475 & 0,034 & 0,063 & 0,392 & 0,007 \\
             & \textit{Skills} + RAG & 0,561 & 0,354 & 0,030 & 0,056 & 0,385 & $-$0,012 \\
             & Completo & 0,565 & 0,378 & 0,025 & 0,047 & 0,383 & $-$0,006 \\
            \bottomrule
        \end{tabular*}
    }
    \endgroup
\end{table}

\begin{table}[H]
    \centering
    \begingroup
    \small
    \setlength{\tabcolsep}{4pt}
    \ttabbox[\FBwidth]{
        \caption{Ablación por componentes en asimetría entre mujeres y hombres (V33), nivel B1 ($N=1.313$, prevalencia 0,062). Los modelos se nombran de forma abreviada, gemini por gemini-3.1-flash-lite y gemma por gemma4:e4b.}
        \label{tab:abl-v33}
    }{
        \begin{tabular*}{\textwidth}{@{\extracolsep{\fill}}l l c c c c c c@{}}
            \toprule
            \textbf{Modelo} & \textbf{Configuración} & \textbf{Exac.} & \textbf{Prec.} & \textbf{Recall} & \textbf{F1 (\textit{Sí})} & \textbf{F1 macro} & \textbf{$\kappa$} \\
            \midrule
            gemini & Solo \textit{skills} & 0,910 & 0,122 & 0,074 & 0,092 & 0,523 & 0,048 \\
             & \textit{Skills} + resúmenes & 0,900 & 0,109 & 0,086 & 0,097 & 0,522 & 0,045 \\
             & \textit{Skills} + RAG & 0,909 & 0,135 & 0,086 & 0,105 & 0,529 & 0,060 \\
             & Completo & 0,909 & 0,118 & 0,074 & 0,091 & 0,521 & 0,045 \\
            \midrule
            gpt-4o-mini & Solo \textit{skills} & 0,796 & 0,061 & 0,160 & 0,088 & 0,487 & $-$0,001 \\
             & \textit{Skills} + resúmenes & 0,781 & 0,081 & 0,247 & 0,122 & 0,498 & 0,032 \\
             & \textit{Skills} + RAG & 0,727 & 0,074 & 0,296 & 0,118 & 0,478 & 0,021 \\
             & Completo & 0,813 & 0,063 & 0,148 & 0,089 & 0,492 & 0,003 \\
            \midrule
            gpt-5.4-nano & Solo \textit{skills} & 0,931 & 0,000 & 0,000 & 0,000 & 0,482 & $-$0,012 \\
             & \textit{Skills} + resúmenes & 0,933 & 0,111 & 0,012 & 0,022 & 0,494 & 0,010 \\
             & \textit{Skills} + RAG & 0,936 & 0,000 & 0,000 & 0,000 & 0,483 & $-$0,004 \\
             & Completo & 0,931 & 0,000 & 0,000 & 0,000 & 0,482 & $-$0,012 \\
            \midrule
            gemma & Solo \textit{skills} & 0,896 & 0,174 & 0,185 & 0,180 & 0,562 & 0,124 \\
             & \textit{Skills} + resúmenes & 0,898 & 0,156 & 0,148 & 0,152 & 0,549 & 0,098 \\
             & \textit{Skills} + RAG & 0,883 & 0,109 & 0,123 & 0,116 & 0,527 & 0,054 \\
             & Completo & 0,889 & 0,143 & 0,160 & 0,151 & 0,546 & 0,092 \\
            \bottomrule
        \end{tabular*}
    }
    \endgroup
\end{table}

\begin{table}[H]
    \centering
    \begingroup
    \small
    \setlength{\tabcolsep}{4pt}
    \ttabbox[\FBwidth]{
        \caption{Ablación por componentes en denominación sexualizada (V35), nivel B1 ($N=1.313$, prevalencia 0,100). Los modelos se nombran de forma abreviada, gemini por gemini-3.1-flash-lite y gemma por gemma4:e4b.}
        \label{tab:abl-v35}
    }{
        \begin{tabular*}{\textwidth}{@{\extracolsep{\fill}}l l c c c c c c@{}}
            \toprule
            \textbf{Modelo} & \textbf{Configuración} & \textbf{Exac.} & \textbf{Prec.} & \textbf{Recall} & \textbf{F1 (\textit{Sí})} & \textbf{F1 macro} & \textbf{$\kappa$} \\
            \midrule
            gemini & Solo \textit{skills} & 0,897 & 0,250 & 0,015 & 0,029 & 0,487 & 0,017 \\
             & \textit{Skills} + resúmenes & 0,901 & 0,571 & 0,031 & 0,058 & 0,503 & 0,048 \\
             & \textit{Skills} + RAG & 0,896 & 0,222 & 0,015 & 0,029 & 0,487 & 0,016 \\
             & Completo & 0,894 & 0,167 & 0,015 & 0,028 & 0,486 & 0,011 \\
            \midrule
            gpt-4o-mini & Solo \textit{skills} & 0,896 & 0,368 & 0,053 & 0,093 & 0,519 & 0,070 \\
             & \textit{Skills} + resúmenes & 0,894 & 0,333 & 0,061 & 0,103 & 0,523 & 0,075 \\
             & \textit{Skills} + RAG & 0,884 & 0,294 & 0,115 & 0,165 & 0,551 & 0,115 \\
             & Completo & 0,897 & 0,375 & 0,046 & 0,082 & 0,514 & 0,061 \\
            \midrule
            gpt-5.4-nano & Solo \textit{skills} & 0,899 & 0,000 & 0,000 & 0,000 & 0,473 & $-$0,003 \\
             & \textit{Skills} + resúmenes & 0,901 & 0,667 & 0,015 & 0,030 & 0,489 & 0,025 \\
             & \textit{Skills} + RAG & 0,896 & 0,222 & 0,015 & 0,029 & 0,487 & 0,016 \\
             & Completo & 0,901 & 1,000 & 0,008 & 0,015 & 0,482 & 0,014 \\
            \midrule
            gemma & Solo \textit{skills} & 0,896 & 0,381 & 0,061 & 0,105 & 0,525 & 0,080 \\
             & \textit{Skills} + resúmenes & 0,891 & 0,300 & 0,069 & 0,112 & 0,527 & 0,078 \\
             & \textit{Skills} + RAG & 0,893 & 0,273 & 0,046 & 0,078 & 0,511 & 0,051 \\
             & Completo & 0,890 & 0,158 & 0,023 & 0,040 & 0,491 & 0,015 \\
            \bottomrule
        \end{tabular*}
    }
    \endgroup
\end{table}


La primera lectura afecta a la \textbf{recuperación en vivo}. La configuración con
RAG rara vez es la mejor de su modelo, y su aportación no se sostiene en ninguna
variable de forma consistente. Dado que es además el componente más costoso en
tiempo y en \textit{tokens}, la evidencia no respalda mantenerlo, conclusión que ya
apuntaba la implementación al describir las limitaciones del índice léxico
(Sección~\ref{rag}).

La segunda lectura es menos cómoda y responde a una pregunta legítima, la de si
procede conservar la configuración completa. \textbf{La configuración completa no
es la mejor en la mayoría de las variables.} Atendiendo en cada una a la métrica
que la Sección~\ref{sec:metricas_eval} considera informativa, el cuadro es el
siguiente. En \texttt{masc\_generico} (V26) la configuración completa de
\texttt{gemini-3.1-flash-lite} obtiene el mejor F1 macro del conjunto (0,621, con
$\kappa = 0,261$), si bien la variante sin RAG queda a dos milésimas (0,619) con un
$\kappa$ ligeramente superior (0,266). En \texttt{sexismo\_discurso} (V30) la mejor
combinación es \texttt{gpt-5.4-nano} con \textit{skills} y RAG (F1 de 0,270 frente
al 0,181 de la completa). En \texttt{asimetria\_mujer\_hombre} (V33), donde importa
no perder positivos, \texttt{gemma4:e4b} con solo \textit{skills} alcanza el mejor
equilibrio (recall 0,185 con precisión 0,174, $\kappa = 0,124$). En
\texttt{denominacion\_sexualizada} (V35) es \texttt{gpt-4o-mini} con \textit{skills}
y RAG (recall 0,115 con precisión 0,294). En \texttt{lenguaje\_sexista} (V25),
finalmente, ninguna configuración alcanza un acuerdo apreciable, ya que todos los
valores de $\kappa$ rondan el cero o son negativos, de modo que la comparación entre
componentes carece ahí de sentido.

De ello se sigue que la arquitectura no admite una respuesta única, sino que
\textbf{la configuración conveniente depende de la variable}. Este resultado
matiza la propuesta inicial, que asumía una configuración uniforme para las cinco,
y sugiere que un despliegue razonable en el proyecto IRIS ajustaría los componentes
variable a variable en lugar de aplicar el sistema completo a todas.

Conviene añadir una advertencia sobre la magnitud de estas diferencias. Las
distancias entre configuraciones son en general pequeñas, en ocasiones de milésimas,
y este trabajo no dispone de intervalos de confianza ni de repeticiones con
distintas semillas que permitan afirmar que sean estadísticamente significativas.
Cabe sostener que la configuración completa no muestra ventaja, pero no que una
configuración concreta sea definitivamente superior a otra.

\subsection{Diagnóstico por variable}
\label{subsec:iris_porvariable}
\label{subsec:iris_diagnostico}

Hasta ahora el análisis se ha apoyado en métricas globales, promediadas sobre las
cinco variables. Ese nivel agregado responde a la pregunta central del capítulo,
pero esconde comportamientos muy distintos de una variable a otra. Para entender qué
ocurre por debajo del promedio hay que desglosar los resultados variable a variable:
primero para localizar dónde actúan las \textit{skills} (B0 frente a B1) y después
para caracterizar el comportamiento del sistema final (B1) en cada variable.

\subsubsection{Efecto de las Agent Skills por variable (B0 frente a B1)}

El desglose por variable revela dónde actúan las \textit{skills}. La
Tabla~\ref{tab:iris-porvariable-b0b1} compara el $\kappa$ de B0 y B1 en cada
variable. En \texttt{gemini}, la mejora del núcleo se concentra por completo en
\texttt{masc\_generico} (de 0,073 a 0,261), mientras el resto de variables apenas se
mueve. En los otros tres modelos ocurre lo contrario: las \textit{skills} degradan
variables donde el control ya tenía algo de señal, sobre todo
\texttt{denominacion\_sexualizada} (de 0,123 a 0,014 en \texttt{gpt-5.4-nano}, de
0,136 a 0,061 en \texttt{gpt-4o-mini} y de 0,154 a 0,015 en \texttt{gemma4:e4b}). El
efecto de la arquitectura no solo depende del modelo, sino que dentro de cada modelo
se localiza en variables concretas.

La Tabla~\ref{tab:iris-porvariable-b0b1-acc} recoge la exactitud en esa misma
comparación, y su movimiento no acompaña al del $\kappa$. En
\texttt{masc\_generico} de \texttt{gemini}, exactitud y $\kappa$ suben a la vez (de
0,397 a 0,708 de exactitud), pero en otras celdas la exactitud cambia en sentido
contrario al acuerdo, como en \texttt{asimetria\_mujer\_hombre} de
\texttt{gpt-5.4-nano}, donde sube de 0,619 a 0,931 mientras el $\kappa$ cae de 0,026
a $-0,012$. Es una nueva manifestación del desacoplamiento de la
Sección~\ref{subsec:iris_divergencia}: la exactitud por variable sigue el desequilibrio,
no el acuerdo.

\begin{table}[H]
    \centering
    \begingroup
    \scriptsize
    \ttabbox[\FBwidth]{
        \caption{Coeficiente $\kappa$ por variable en B0 (\textit{baseline}) y B1 (Agent Skills), para los cuatro modelos ($N=1.313$). Cada celda muestra B0 / B1.}
        \label{tab:iris-porvariable-b0b1}
    }{
        \begin{tabularx}{\textwidth}{
            >{\hsize=1.4\hsize\raggedright\arraybackslash\ttfamily}X
            >{\hsize=0.9\hsize\centering\arraybackslash}X
            >{\hsize=0.9\hsize\centering\arraybackslash}X
            >{\hsize=0.9\hsize\centering\arraybackslash}X
            >{\hsize=0.9\hsize\centering\arraybackslash}X
        }
            \toprule
            \normalfont\textbf{Variable} & \textbf{Gemini} & \textbf{GPT-4o-mini} & \textbf{GPT-5.4-nano} & \textbf{Gemma} \\
            \midrule
            lenguaje\_sexista & 0,000 / 0,015 & $-0,013$ / $-0,037$ & 0,001 / 0,000 & $-0,004$ / 0,003 \\
            masc\_generico & \textbf{0,073 / 0,261} & 0,023 / 0,034 & 0,103 / 0,016 & 0,102 / 0,109 \\
            sexismo\_discurso & 0,004 / 0,000 & 0,002 / 0,011 & 0,007 / 0,004 & 0,002 / $-0,006$ \\
            asimetria\_mujer\_hombre & 0,036 / 0,045 & 0,019 / 0,003 & 0,026 / $-0,012$ & 0,054 / 0,092 \\
            denominacion\_sexualizada & 0,019 / 0,011 & 0,136 / 0,061 & 0,123 / 0,014 & 0,154 / 0,015 \\
            \bottomrule
        \end{tabularx}
    }
    \endgroup
\end{table}

\begin{table}[H]
    \centering
    \begingroup
    \scriptsize
    \ttabbox[\FBwidth]{
        \caption{Exactitud por variable en B0 (\textit{baseline}) y B1 (Agent Skills), para los cuatro modelos ($N=1.313$). Cada celda muestra B0 / B1.}
        \label{tab:iris-porvariable-b0b1-acc}
    }{
        \begin{tabularx}{\textwidth}{
            >{\hsize=1.4\hsize\raggedright\arraybackslash\ttfamily}X
            >{\hsize=0.9\hsize\centering\arraybackslash}X
            >{\hsize=0.9\hsize\centering\arraybackslash}X
            >{\hsize=0.9\hsize\centering\arraybackslash}X
            >{\hsize=0.9\hsize\centering\arraybackslash}X
        }
            \toprule
            \normalfont\textbf{Variable} & \textbf{Gemini} & \textbf{GPT-4o-mini} & \textbf{GPT-5.4-nano} & \textbf{Gemma} \\
            \midrule
            lenguaje\_sexista & 0,200 / 0,226 & 0,236 / 0,457 & 0,207 / 0,191 & 0,199 / 0,201 \\
            masc\_generico & \textbf{0,397 / 0,708} & 0,317 / 0,332 & 0,580 / 0,262 & 0,451 / 0,482 \\
            sexismo\_discurso & 0,571 / 0,569 & 0,554 / 0,561 & 0,554 / 0,558 & 0,567 / 0,565 \\
            asimetria\_mujer\_hombre & 0,936 / 0,909 & 0,816 / 0,813 & 0,619 / 0,931 & 0,877 / 0,889 \\
            denominacion\_sexualizada & 0,898 / 0,894 & 0,896 / 0,897 & 0,895 / 0,901 & 0,893 / 0,890 \\
            \bottomrule
        \end{tabularx}
    }
    \endgroup
\end{table}

\subsubsection{Comportamiento por variable en el nivel B1}

Para cada una se reportan, sobre la clase positiva (\textit{Sí}), la \textbf{precisión} (de
las piezas que el modelo marca como \textit{positivas}, cuántas lo son de verdad) y el
\textbf{recall} o sensibilidad (del sexismo que marcan las expertas, cuánto detecta el
modelo), junto a la \textbf{exactitud} (el acierto global sobre todas las piezas, que
no debe confundirse con la precisión) y el coeficiente $\kappa$. El recall es aquí el
más revelador: un recall bajo significa que el modelo deja sin marcar buena parte de lo
que las anotadoras sí marcan. La Tabla~\ref{tab:iris-diagnostico} recoge estos valores
para el nivel B1.


\begin{table}[H]
    \centering
    \begingroup
    \scriptsize
    \setlength{\tabcolsep}{3pt}
    \newcommand{\cel}[4]{\begin{tabular}[t]{@{}c@{}}#1\\#2\\#3\\#4\end{tabular}}
    \ttabbox[\FBwidth]{
        \caption{Exactitud, precisión, recall (sensibilidad) sobre la clase \textit{Sí} y $\kappa$, por variable en el nivel B1 ($N=1.313$). En cada celda se apilan, de arriba abajo, exactitud, precisión, recall y $\kappa$.}
        \label{tab:iris-diagnostico}
    }{
        \begin{tabularx}{\textwidth}{
            >{\hsize=1.35\hsize\raggedright\arraybackslash\ttfamily}X
            >{\hsize=0.55\hsize\centering\arraybackslash}X
            >{\hsize=0.775\hsize\centering\arraybackslash}X
            >{\hsize=0.775\hsize\centering\arraybackslash}X
            >{\hsize=0.775\hsize\centering\arraybackslash}X
            >{\hsize=0.775\hsize\centering\arraybackslash}X
        }
            \toprule
            \normalfont\textbf{Variable} & \textbf{Prev. Sí} & \textbf{Gemini} & \textbf{GPT-4o-mini} & \textbf{GPT-5.4-nano} & \textbf{Gemma} \\
            \midrule
            lenguaje\_sexista & 0,810 & \cel{0,226}{0,961}{0,046}{0,015} & \cel{0,457}{0,789}{0,450}{$-$0,037} & \cel{0,191}{1,000}{0,001}{0,000} & \cel{0,201}{0,929}{0,012}{0,003} \\
            \midrule
            masc\_generico & 0,810 & \textbf{\cel{0,708}{0,887}{0,733}{0,261}} & \cel{0,332}{0,873}{0,206}{0,034} & \cel{0,262}{0,874}{0,104}{0,016} & \cel{0,482}{0,897}{0,408}{0,109} \\
            \midrule
            sexismo\_discurso & 0,428 & \cel{0,569}{0,423}{0,020}{0,000} & \cel{0,561}{0,450}{0,112}{0,011} & \cel{0,558}{0,435}{0,114}{0,004} & \cel{0,565}{0,378}{0,025}{$-$0,006} \\
            \midrule
            asimetria\_mujer\_hombre & 0,062 & \cel{0,909}{0,118}{0,074}{0,045} & \cel{0,813}{0,063}{0,148}{0,003} & \cel{0,931}{0,000}{0,000}{$-$0,012} & \cel{0,889}{0,143}{0,160}{0,092} \\
            \midrule
            denominacion\_sexualizada & 0,100 & \cel{0,894}{0,167}{0,015}{0,011} & \cel{0,897}{0,375}{0,046}{0,061} & \cel{0,901}{1,000}{0,008}{0,014} & \cel{0,890}{0,158}{0,023}{0,015} \\
            \bottomrule
        \end{tabularx}
    }
    \endgroup
\end{table}


\begin{itemize}
    \item \textbf{\texttt{lenguaje\_sexista} (V25): detección fallida.} Es el
    resultado más severo del capítulo. El recall es prácticamente nulo en tres de los
    cuatro modelos (0,046 en \texttt{gemini}, 0,012 en \texttt{gemma} y 0,001 en
    \texttt{gpt-5.4-nano}), de modo que el sistema deja sin señalar casi todo el
    lenguaje sexista que las expertas identifican. No se trata de un matiz de umbral
    ni de un desacuerdo sobre casos dudosos, sino de que \textbf{el modelo no detecta
    el fenómeno}. Y conviene subrayar la magnitud de lo que se está pasando por alto,
    ya que ambos equipos de anotación marcan esta variable en la mayoría de las
    piezas, el 74,6\,\% Indexa y el 85,5\,\% UCM3
    (Tabla~\ref{tab:iris-equipos-var}), es decir, coinciden en que el fenómeno está
    ahí aunque discrepen en cuánto. La única excepción es \texttt{gpt-4o-mini}, que sí
    marca de forma abundante (recall 0,450), pero su acuerdo sigue siendo negativo
    ($\kappa = -0,037$), lo que indica que marca en piezas distintas de aquellas en
    que lo hacen las anotadoras. Ni detectar poco ni detectar mucho aproxima al modelo
    al criterio experto. Que la variable sea la más genérica del bloque, el uso
    discriminatorio del lenguaje por razón de sexo sin describir un mecanismo
    concreto, ayuda a explicar la dificultad, pero no la disuelve: identificar ese uso
    es precisamente lo que se pedía al sistema. Queda por tanto un trabajo
    considerable por delante antes de que un modelo pueda asistir en esta variable, y
    se retoma en la discusión (Sección~\ref{sec:discussion}).

    \item \textbf{\texttt{masc\_generico} (V26): el único acuerdo genuino, y su
    umbral.} Es el único caso con acuerdo apreciable: \texttt{gemini} alcanza
    $\kappa = 0,261$ y un recall de 0,733, es decir, detecta la mayoría del masculino
    genérico que marcan las expertas, con exactitud y $\kappa$ que mejoran a la vez
    sobre el \textit{baseline} (lo que descarta el artefacto de marginales y confirma
    una ganancia real). El local \texttt{gemma4:e4b} queda en segundo lugar, con un
    recall de 0,408 y un $\kappa$ de 0,109, la única variable en la que también supera
    con claridad el azar. Los dos modelos de OpenAI, en cambio, son mucho más
    conservadores (recall 0,206 y 0,104): marcan el masculino genérico en muy pocas
    piezas. La variable es propensa a este desajuste de umbral, porque el español lo
    emplea de forma pervasiva (\textit{los trabajadores}, \textit{los ciudadanos}) y la
    anotación tiende a marcarlo de forma sistemática, de modo que un modelo restrictivo
    se queda corto. El umbral en el que cada modelo decide marcar es una característica
    del propio modelo, y los de OpenAI lo sitúan más alto que \texttt{gemini} o
    \texttt{gemma}.

    \item \textbf{\texttt{sexismo\_discurso} (V30): baja detección, pero por otro
    motivo.} El acuerdo ronda el azar en los cuatro modelos ($\kappa$ entre $-0,006$ y
    0,011) y el recall es muy bajo (entre 0,020 y 0,114), de modo que, pese a una
    prevalencia anotada del 42,8\,\%, los modelos apenas marcan esta variable. La
    infradetección se parece a la de V25, pero su origen no es el mismo. Aquí la
    propia anotación humana está profundamente dividida, ya que un equipo marca el
    fenómeno casi seis veces más que el otro (11,5\,\% frente a 65,2\,\%,
    Tabla~\ref{tab:iris-equipos-var}), mientras que en V25 ambos equipos coincidían en
    marcar la mayoría de las piezas. En V25 los modelos no ven lo que las personas ven
    de forma concordante. En V30, en cambio, ni siquiera las personas convergen en qué
    debe marcarse, así que el bajo acuerdo del sistema refleja a la vez una detección
    insuficiente y un criterio humano poco estabilizado. Es además la única variable de
    \emph{fondo} de las cinco, la que atañe a lo que se cuenta y no a cómo se dice
    (Sección~\ref{sec:iris_variables}), lo que la hace más dependiente de la
    interpretación del contexto social de la pieza.

    \item \textbf{\texttt{asimetria\_mujer\_hombre} (V33) y
    \texttt{denominacion\_sexualizada} (V35): desequilibrio extremo.} Con prevalencias
    del 6,2\,\% y el 10,0\,\%, son las dos variables en que la exactitud resulta menos
    informativa y en que el F1 de la clase positiva, concebido precisamente para
    situaciones de clase minoritaria, ofrece la lectura más ajustada
    (Sección~\ref{sec:metricas_eval}). Los valores son bajos (entre 0,000 y 0,151),
    con recalls que no superan 0,160, de manera que también aquí el sistema deja sin
    señalar la mayor parte de los casos. Merece un matiz el comportamiento de
    \texttt{gpt-5.4-nano} en \texttt{denominacion\_sexualizada}, que marca \textit{Sí}
    en tan pocas piezas que su precisión de 1,000 descansa sobre un puñado de
    predicciones y no permite conclusión alguna sobre su fiabilidad general. Dicho
    esto, y aunque un recall de 0,008 lo vuelve inservible como detector, el dato
    tiene una lectura menos desfavorable, y es que cuando ese modelo decide marcar,
    acierta. Un sistema así resulta inútil para revisar un corpus completo, pero no
    genera falsas alarmas.

\end{itemize}

\newpage

Conviene mirar también cuánto marca cada modelo en términos absolutos, que ayuda a
interpretar el recall: marcar poco impide detectar buena parte de lo que las
expertas sí señalan. La Tabla~\ref{tab:iris-prev-modelo-var} recoge, por variable, la
prevalencia de \textit{Sí} de cada modelo. La única variable en la que algún modelo
se acerca al terreno de las expertas es \texttt{masc\_generico} (\texttt{gemini},
66,9\,\%), mientras que en el resto todos marcan \textit{Sí} en una fracción muy
pequeña de las piezas. El patrón, además, no es común: \texttt{gpt-4o-mini} dispara
en \texttt{lenguaje\_sexista} (44,6\,\%) donde \texttt{gemini} apenas marca (3,8\,\%),
mientras que \texttt{gemini} domina en \texttt{masc\_generico}. Cada modelo tiene su
propio umbral por variable.

\begin{table}[H]
    \centering
    \begingroup
    \scriptsize
    \renewcommand{\texttt}[1]{{\ttfamily #1}}
    \ttabbox[\FBwidth]{
        \caption{Prevalencia de \textit{Sí} (\%) de cada modelo por variable, en el nivel B1 ($N=1.313$). La columna \textbf{Prev. Sí} (sombreada) es la prevalencia anotada por las expertas, que sirve de referencia.}
        \label{tab:iris-prev-modelo-var}
    }{
        \begin{tabularx}{\textwidth}{
            >{\hsize=2\hsize\raggedright\arraybackslash\ttfamily}X
            >{\hsize=0.7875\hsize\columncolor{gray!12}\centering\arraybackslash}X
            >{\hsize=0.7875\hsize\centering\arraybackslash}X
            >{\hsize=0.7875\hsize\centering\arraybackslash}X
            >{\hsize=0.7875\hsize\centering\arraybackslash}X
            >{\hsize=0.7875\hsize\centering\arraybackslash}X
        }
            \toprule
            \normalfont\textbf{Variable} & \textbf{Prev. Sí} & \textbf{Gemini} & \textbf{GPT-4o-mini} & \textbf{GPT-5.4-nano} & \textbf{Gemma} \\
            \midrule
            lenguaje\_sexista & 81,0 & 3,8 & 44,6 & 0,1 & 1,1 \\
            masc\_generico & 81,0 & \textbf{66,9} & 19,1 & 9,7 & 36,9 \\
            sexismo\_discurso & 42,8 & 2,0 & 10,7 & 11,2 & 2,8 \\
            asimetria\_mujer\_hombre & 6,2 & 3,9 & 14,4 & 0,7 & 6,9 \\
            denominacion\_sexualizada & 10,0 & 0,9 & 1,2 & 0,1 & 1,4 \\
            \bottomrule
        \end{tabularx}
    }
    \endgroup
\end{table}

Un rasgo atraviesa a las cinco variables: el recall es bajo casi en todas, es decir,
los modelos detectan solo una fracción del sexismo que marcan las expertas. Esta
\textbf{infradetección generalizada} es el patrón dominante del capítulo. La
Sección~\ref{subsec:iris_errores} la examina para determinar \emph{qué parte} de ese
desajuste corresponde a una limitación de los modelos y qué parte a la dificultad de
la propia tarea, sin presuponer que se deba únicamente a una de las dos. A ello se
añade el comportamiento dispar de los modelos según la variable, que desaconseja
tratarlos como un bloque homogéneo.



\subsection{El modelo local: rendimiento y coste}
\label{subsec:iris_local}\label{subsec:iris_coste}

Entre los cuatro modelos, el local \texttt{gemma4:e4b} merece una lectura aparte. En
el nivel de control (B0) obtiene el \textbf{mayor $\kappa$ de los cuatro} (0,061), por
encima de las tres API de pago (0,026, 0,033 y 0,052), y ante las \textit{skills} se
comporta como los modelos de OpenAI, ya que su acuerdo medio baja de 0,061 a
0,043. Que un modelo abierto, gratuito y desplegable
en infraestructura propia iguale, y en $\kappa$ supere, a los servicios
comerciales tiene consecuencias prácticas para un proyecto que maneja material
sujeto a privacidad. Ahora bien, esa ventaja en acuerdo y en ausencia de tarifas
convive con un coste de otro tipo: el del cómputo.

La Tabla~\ref{tab:iris-b0b1} recoge el coste real de inferencia. Sobre las 1.313
noticias, el sistema agéntico (B1) cuesta 30,91 USD en los tres modelos de pago,
frente a los 13,83 USD del \textit{baseline} (B0), un factor de 2,2. El aumento se
explica por el número de llamadas al modelo: B0 resuelve cada variable en una única
llamada, mientras que en B1 la agente encadena varias (entre dos y tres por variable
en promedio, las acciones de su bucle más el veredicto final). Que el factor no sea
aún mayor se debe a que cada llamada de B1 es más corta: gracias a la carga bajo
demanda, la agente solo incluye en el \textit{prompt} lo que necesita, mientras que
B0 inyecta la metodología completa en su única llamada.


En el modelo local no hay tarifas, pero el coste reaparece como tiempo de cómputo. El
nivel B1 de \texttt{gemma4:e4b} se ejecutó sobre la granja de GPU (NVIDIA RTX 3090 y
4090) del Departamento de Teoría de la Señal y Comunicaciones (TSC) de la Universidad
Carlos III de Madrid, con un rendimiento medido de unos 44 segundos por variable,
esto es, en torno a 220 segundos por artículo completo. Repartiendo la carga en
cuatro servidores en paralelo (uno por \textit{shard}), cada ejecución sobre el corpus
de evaluación tardó alrededor de veinte horas. En un solo servidor, el mismo cómputo
superaría las ochenta horas. Así, el modelo
local es competitivo en acuerdo y prescinde de tarifas, pero su despliegue agéntico
exige infraestructura dedicada: el coste no desaparece, cambia de dinero a cómputo.
La Figura~\ref{fig:iris-acuerdo-dolar} resume ese compromiso: situado el acuerdo del
sistema desplegado (B1) frente a su coste monetario, el modelo local ocupa la esquina
más favorable (pintada en verde). Sin ser el de mayor acuerdo absoluto (\texttt{gemini} lo supera en
$\kappa$), ofrece un acuerdo competitivo a coste monetario nulo, con la salvedad de
que ese coste se traslada al cómputo. Las implicaciones prácticas para IRIS se
retoman en la discusión (Sección~\ref{sec:discussion}).

\begin{figure}[H]
    \centering
    \includegraphics[width=\textwidth]{fig_acuerdo_por_dolar_blanco.pdf}
    \caption{Acuerdo (coeficiente $\kappa$ en B1) frente al coste de inferencia sobre
    el corpus completo, por modelo ($N=1.313$). El modelo local ofrece el mejor
    acuerdo por dólar (coste monetario nulo). Su coste se traslada al tiempo de
    cómputo en infraestructura propia.}
    \label{fig:iris-acuerdo-dolar}
\end{figure}

\subsection{¿Un límite del modelo o de la tarea?}
\label{subsec:iris_errores}

El bajo acuerdo absoluto plantea la pregunta que condiciona la interpretación de
todo el capítulo: ¿refleja una limitación de los modelos, o un límite intrínseco de
la propia tarea? Tres análisis complementarios apuntan en la misma dirección.

\subsubsection{Heterogeneidad entre equipos de anotación}

El corpus registra, para cada pieza, la persona que la codificó. Las 1.313
noticias anotadas se reparten entre dos equipos: Indexa (548 piezas) y UCM3
(765). Aunque la ausencia de doble codificación
(Sección~\ref{subsec:limitacion_fiabilidad}) impide calcular la fiabilidad
intercodificadora, sí es posible medir el acuerdo del sistema con cada equipo
por separado, a modo de análisis de sensibilidad.

\begin{table}[H]
    \centering
    \begingroup
    \scriptsize
    \renewcommand{\texttt}[1]{{\ttfamily #1}}
    \ttabbox[\FBwidth]{
        \caption{Análisis de sensibilidad: acuerdo (B1) según el equipo de anotación}
        \label{tab:iris-sensibilidad}
    }{
        \begin{tabularx}{\textwidth}{
            >{\hsize=1.5\hsize\raggedright\arraybackslash}X
            >{\hsize=1.0\hsize\raggedright\arraybackslash}X
            >{\hsize=0.85\hsize\centering\arraybackslash}X
            >{\hsize=0.85\hsize\centering\arraybackslash}X
            >{\hsize=0.8\hsize\columncolor{gray!15}\centering\arraybackslash}X
        }
            \toprule
            \textbf{Modelo} & \textbf{Subconjunto} & \textbf{Exactitud} & \textbf{F1 macro} & \boldmath$\kappa$ \\
            \midrule
            \multirow{3}{*}{\texttt{gemini-3.1-flash-lite}}
              & Todas (1.313)   & 0,661 & 0,443 & 0,066 \\
              & \textbf{Solo Indexa (548)} & \textbf{0,741} & \textbf{0,503} & \textbf{0,118} \\
              & Solo UCM3 (765)  & 0,604 & 0,398 & 0,043 \\
            \midrule
            \multirow{3}{*}{\texttt{gpt-4o-mini}}
              & Todas           & 0,612 & 0,438 & 0,014 \\
              & \textbf{Solo Indexa}     & \textbf{0,693} & \textbf{0,482} & \textbf{0,042} \\
              & Solo UCM3       & 0,554 & 0,398 & 0,001 \\
            \midrule
            \multirow{3}{*}{\texttt{gpt-5.4-nano}}
              & Todas           & 0,569 & 0,364 & 0,004 \\
              & \textbf{Solo Indexa}     & \textbf{0,661} & \textbf{0,406} & \textbf{0,014} \\
              & Solo UCM3       & 0,502 & 0,323 & $-0,001$ \\
            \midrule
            \multirow{3}{*}{\texttt{gemma4:e4b}}
              & Todas           & 0,605 & 0,412 & 0,043 \\
              & \textbf{Solo Indexa}     & \textbf{0,695} & \textbf{0,461} & \textbf{0,068} \\
              & Solo UCM3       & 0,541 & 0,369 & 0,032 \\
            \bottomrule
        \end{tabularx}
    }
    \endgroup
\end{table}

La Tabla~\ref{tab:iris-sensibilidad} muestra un resultado nítido: el acuerdo con
Indexa es sistemáticamente superior al obtenido con UCM3, en los cuatro modelos y en
las tres métricas. En \texttt{gemini-3.1-flash-lite}, restringir la evaluación a
Indexa casi duplica el coeficiente $\kappa$ (de 0,066 a 0,118) y sube la exactitud
ocho puntos, mientras que contra UCM3 el $\kappa$ cae hacia el azar (llega a valores
negativos en \texttt{gpt-5.4-nano}).

Lo interesante es la magnitud de ese efecto: \textbf{cambiar de equipo de anotación
mueve el acuerdo tanto o más que cambiar de modelo.} En \texttt{gemini}, el salto de
$\kappa$ entre evaluar contra Indexa y contra UCM3 (de 0,118 a 0,043, un rango de
0,075, Tabla~\ref{tab:iris-sensibilidad}) supera al rango completo de $\kappa$ entre
los cuatro modelos, que va de 0,004 en \texttt{gpt-5.4-nano} a 0,066 en el propio
\texttt{gemini} (un rango de 0,062, Tabla~\ref{tab:iris-b0b1}). Es decir, contra qué
equipo se compara pesa en el resultado tanto como qué modelo se emplea.


El origen de esa diferencia está en la anotación, no en el modelo. La
Tabla~\ref{tab:iris-equipos-var} muestra, por variable, la proporción de \textit{Sí}
que marca cada equipo. Esas proporciones son de la propia anotación humana y, por
tanto, \textbf{no dependen de ningún modelo}: reflejan cómo codifica cada equipo, y
dejan ver que la misma variable se entiende de forma muy distinta. El caso extremo es
\texttt{sexismo\_discurso}: Indexa marca el fenómeno en el 11,5\,\% de las piezas y
UCM3 en el 65,2\,\%, casi seis veces más. No es ruido del modelo, son dos criterios
humanos difícilmente conciliables sobre una misma variable. Las dos últimas columnas
de la tabla añaden, a modo de ilustración, el $\kappa$ de \texttt{gemini} frente a
cada equipo. El mismo patrón se da en los demás modelos, como ya recogía el análisis
de sensibilidad (Tabla~\ref{tab:iris-sensibilidad}).
Como esa variabilidad entre equipos iguala o supera la que separa a los modelos,
constituye una evidencia del techo de la tarea, sin necesitar doble codificación.


\begin{table}[H]
    \centering
    \begingroup
    \scriptsize
    \renewcommand{\texttt}[1]{{\ttfamily #1}}
    \ttabbox[\FBwidth]{
        \caption{Prevalencia de \textit{Sí} y acuerdo (\texttt{gemini} B1) por equipo y variable. La misma variable se codifica de forma muy distinta según el equipo.}
        \label{tab:iris-equipos-var}
    }{
        \begin{tabularx}{\textwidth}{
            >{\hsize=1.7\hsize\raggedright\arraybackslash\ttfamily}X
            >{\hsize=0.9\hsize\centering\arraybackslash}X
            >{\hsize=0.9\hsize\centering\arraybackslash}X
            >{\hsize=0.9\hsize\centering\arraybackslash}X
            >{\hsize=0.8\hsize\centering\arraybackslash}X
            >{\hsize=0.8\hsize\centering\arraybackslash}X
        }
            \toprule
             & \multicolumn{2}{c}{\normalfont\% \textit{Sí}} & & \multicolumn{2}{c}{\normalfont$\kappa$} \\
            \cmidrule(lr){2-3}\cmidrule(lr){5-6}
            \normalfont\textbf{Variable} & \textbf{Indexa} & \textbf{UCM3} & \textbf{Ratio} & \textbf{Indexa} & \textbf{UCM3} \\
            \midrule
            lenguaje\_sexista & 74,6 & 85,5 & 1,1$\times$ & 0,026 & 0,009 \\
            masc\_generico & 71,4 & 88,0 & 1,2$\times$ & 0,277 & 0,223 \\
            sexismo\_discurso & \textbf{11,5} & \textbf{65,2} & \textbf{5,7$\times$} & 0,065 & $-0,003$ \\
            asimetria\_mujer\_hombre & 3,3 & 8,2 & 2,5$\times$ & 0,202 & $-0,027$ \\
            denominacion\_sexualizada & 6,6 & 12,4 & 1,9$\times$ & 0,019 & 0,013 \\
            \bottomrule
        \end{tabularx}
    }
    \endgroup
\end{table}

Conviene, no obstante, acotar el alcance de este análisis. La mayor concordancia con
Indexa no justifica reportar únicamente ese subconjunto como resultado del sistema:
en ausencia de doble codificación no es posible determinar qué equipo se ajusta mejor
al constructo, y la diferencia observada entre ambos no depende solo de su criterio,
ya que cada equipo codificó piezas distintas. El análisis delimita el techo de acuerdo
alcanzable, pero no lo imputa a un error de anotación de ninguno de los dos equipos.


\subsubsection{Consenso de los modelos frente al \textit{ground truth}}

Un segundo análisis examina las celdas en las que \textbf{los cuatro modelos
coinciden entre sí} pero difieren de la anotación experta. Son celdas informativas,
porque descartan que el desajuste se deba a la particularidad de un modelo concreto,
si bien conviene no sobreinterpretarlas, ya que la coincidencia de los cuatro no
convierte a la anotación en errónea ni demuestra por sí sola una incapacidad general
de los modelos de lenguaje. De
las 6.565 celdas evaluadas (1.313 piezas $\times$ 5 variables), 1.364 presentan
consenso de los cuatro modelos en contra del \textit{ground truth}. Su dirección
es prácticamente unánime: en \textbf{1.363 casos} los cuatro modelos responden \textit{No}
donde la anotación dice \textit{Sí}, y solo en 1 caso ocurre lo contrario. La
infradetección observada por variable es, por tanto, un \textbf{sesgo sistemático},
no aleatorio, y constituye el hallazgo principal de este análisis: \textbf{los
modelos detectan de forma consistente menos sexismo que las anotadoras}.

Este resultado no es aislado. La literatura sobre detección automática de discurso
discriminatorio viene documentando que los sistemas rinden bastante peor cuando el
fenómeno no se apoya en marcas explícitas. ElSherief et al.~\cite{elsherief2021}
muestran que los clasificadores del estado del arte dejan sin marcar buena parte del
odio implícito, con proporciones de casos clasificados como inocuos que llegan al
36\,\% en algunas categorías. La evaluación funcional de
Röttger et al.~\cite{rottger2021} ofrece un dato aún más severo sobre sistemas en
producción, ya que uno de los servicios comerciales examinados solo identifica como
odio el 9,0\,\% de los casos que sí lo son. En español y sobre sexismo,
De Grazia et al.~\cite{degrazia2025} constatan que los modelos de lenguaje resuelven
razonablemente el sexismo explícito pero se atascan en el implícito, que es
precisamente donde también desciende el acuerdo entre anotadoras.

Conviene, no obstante, no elevar esta coincidencia a ley general. Otros trabajos
observan el sesgo contrario, con modelos que marcan en exceso enunciados inocuos que
mencionan a grupos sensibles~\cite{zhang2024extremes}, de modo que la dirección del
error depende del sistema, del umbral y del dominio. Lo que la literatura sí
sostiene con solidez es que el rendimiento cae en lo implícito, y en ese marco los
resultados de este trabajo aportan evidencia propia sobre un dominio poco explorado,
el de la prensa escrita en español evaluada con criterio experto.

El desglose por variable de esas celdas (Tabla~\ref{tab:iris-consenso}) muestra
además \emph{dónde} se concentra el sesgo. Casi tres cuartas partes de los casos se
acumulan en \texttt{lenguaje\_sexista} (567) y \texttt{sexismo\_discurso} (447),
justo las dos variables con mayor prevalencia anotada de \textit{Sí}: donde las
expertas marcan más, los cuatro modelos coinciden en no marcar. La dirección es
prácticamente absoluta, un único caso, en \texttt{asimetria\_mujer\_hombre}, invierte
el patrón (los cuatro modelos marcan \textit{Sí} donde la anotación dice \textit{No}).
Esto refuerza el mensaje del apartado: el sesgo no solo es sistemático (dirección
unánime), sino que se concentra en las variables de alta prevalencia, enlazando con
el diagnóstico por variable (Sección~\ref{subsec:iris_diagnostico}).

\begin{table}[H]
    \centering
    \begingroup
    \scriptsize
    \renewcommand{\texttt}[1]{{\ttfamily #1}}
    \ttabbox[\FBwidth]{
        \caption{Desglose por variable de las 1.364 celdas en que los cuatro modelos coinciden entre sí pero difieren del \textit{ground truth} (B1). ``No/Sí'' = los modelos marcan \textit{No} donde la anotación dice \textit{Sí} (infradetección). ``Sí/No'' = lo contrario (sobredetección).}
        \label{tab:iris-consenso}
    }{
        \begin{tabularx}{\textwidth}{
            >{\hsize=1.9\hsize\raggedright\arraybackslash\ttfamily}X
            >{\hsize=0.9\hsize\centering\arraybackslash}X
            >{\hsize=0.7\hsize\centering\arraybackslash}X
            >{\hsize=0.7\hsize\centering\arraybackslash}X
        }
            \toprule
            \normalfont\textbf{Variable} & \textbf{Consenso $\neq$ GT} & \textbf{No/Sí} & \textbf{Sí/No} \\
            \midrule
            lenguaje\_sexista & 567 & 567 & 0 \\
            masc\_generico & 169 & 169 & 0 \\
            sexismo\_discurso & 447 & 447 & 0 \\
            asimetria\_mujer\_hombre & 61 & 60 & 1 \\
            denominacion\_sexualizada & 120 & 120 & 0 \\
            \midrule
            \textbf{Total} & \textbf{1.364} & \textbf{1.363} & \textbf{1} \\
            \bottomrule
        \end{tabularx}
    }
    \endgroup
\end{table}

\subsubsection{El sistema como una anotadora más}
\label{subsec:iris_equipo_ia}

Los análisis anteriores comparan cada modelo contra la anotación experta tomada
como referencia. Un enfoque complementario invierte la pregunta: en vez de medir
cuánto se equivoca el modelo respecto al \textit{ground truth}, lo sitúa como \textbf{una
anotadora más} dentro del espacio de criterios de las personas que codificaron el
corpus. La cuestión deja de ser \textit{``¿acierta o no el modelo?''} y pasa a ser \textit{``¿dónde se sitúa su
habilidad de detección respecto al de las expertas?''}. Los modelos se tratan así formalmente como un tercer
equipo de anotación, junto a Indexa y UCM3.

Como cada noticia la codificó una sola persona, no hay solapamiento entre anotadoras. Lo que
sí compara todos los criterios es la \textbf{prevalencia del \textit{Sí}}: la proporción de
piezas en que cada entidad codificadora, humana o artificial, marca la variable. La
Tabla~\ref{tab:iris-criterios} ordena por esa proporción, promediada sobre las cinco
variables, a las nueve personas anotadoras y a los cuatro modelos evaluados.

\begin{table}[H]
    \centering
    \begingroup
    \scriptsize
    \renewcommand{\texttt}[1]{{\ttfamily #1}}
    \ttabbox[\FBwidth]{
        \caption{Espacio de criterios: prevalencia de \textit{Sí} por entidad anotadora (humana y artificial) y variable. Columnas: V25 \texttt{lenguaje\_sexista}, V26 \texttt{masc\_generico}, V30 \texttt{sexismo\_discurso}, V33 \texttt{asimetria\_mujer\_hombre}, V35 \texttt{denominacion\_sexualizada}. La última columna (sombreada) promedia las cinco variables y se ofrece solo como
        ordenación de la tabla, ya que con prevalencias tan dispares esa media no
        admite una lectura sustantiva.}
        \label{tab:iris-criterios}
    }{
        \begin{tabularx}{\textwidth}{
            >{\hsize=1.5\hsize\raggedright\arraybackslash}X
            >{\hsize=0.6\hsize\centering\arraybackslash}X
            >{\hsize=0.8\hsize\centering\arraybackslash}X
            >{\hsize=0.8\hsize\centering\arraybackslash}X
            >{\hsize=0.8\hsize\centering\arraybackslash}X
            >{\hsize=0.8\hsize\centering\arraybackslash}X
            >{\hsize=0.8\hsize\centering\arraybackslash}X
            >{\hsize=0.9\hsize\columncolor{gray!15}\centering\arraybackslash}X
        }
            \toprule
            \textbf{Anotador/a} & \boldmath$n$ & \textbf{V25} & \textbf{V26} & \textbf{V30} & \textbf{V33} & \textbf{V35} & \textbf{Media} \\
            \midrule
            UCM3 3 & 363 & 0,986 & 0,975 & 0,978 & 0,041 & 0,215 & 0,639 \\
            UCM3 4 & 166 & 0,831 & 0,855 & 0,283 & 0,271 & 0,084 & 0,465 \\
            Indexa 1 & 127 & 0,772 & 0,756 & 0,307 & 0,087 & 0,118 & 0,408 \\
            UCM3 6 & 1 & 1,000 & 0,000 & 1,000 & 0,000 & 0,000 & 0,400 \\
            Indexa 5 & 50 & 0,940 & 0,920 & 0,100 & 0,020 & 0,000 & 0,396 \\
            UCM3 5 & 235 & 0,668 & 0,753 & 0,409 & 0,013 & 0,013 & 0,371 \\
            Indexa 2 & 128 & 0,742 & 0,727 & 0,062 & 0,031 & 0,062 & 0,325 \\
            Indexa 3 & 231 & 0,701 & 0,649 & 0,039 & 0,009 & 0,048 & 0,289 \\
            Indexa 4 & 12 & 0,583 & 0,500 & 0,167 & 0,000 & 0,167 & 0,283 \\
            \midrule[\heavyrulewidth]
            \texttt{gpt-4o-mini} & 1.313 & 0,462 & 0,191 & 0,107 & 0,144 & 0,012 & 0,183 \\
            \texttt{gemini-3.1-flash-lite} & 1.313 & 0,039 & 0,669 & 0,020 & 0,039 & 0,009 & 0,155 \\
            \texttt{gemma4:e4b} & 1.313 & 0,012 & 0,369 & 0,028 & 0,069 & 0,014 & 0,099 \\
            \texttt{gpt-5.4-nano} & 1.313 & 0,001 & 0,097 & 0,112 & 0,007 & 0,001 & 0,043 \\
            \bottomrule
        \end{tabularx}
    }
    \endgroup
\end{table}

Los cuatro modelos quedan por debajo de la anotadora humana que menos marca, esto
es, fuera del rango humano completo. La diferencia no debe leerse como una postura
más permisiva ante el sexismo, puesto que los modelos no sostienen un criterio
alternativo, sino como una capacidad de detección menor. Es la infradetección
sistemática de la subsección anterior, leída ahora sobre una escala común a personas
y máquinas.

Ese retrato, no obstante, es el del promedio. El desglose por variable matiza el
cuadro, porque en algunas variables ciertos modelos son más inclusivos que la mayoría
de las personas. El caso más claro es \texttt{gpt-4o-mini} en
\texttt{asimetria\_mujer\_hombre}, que marca el fenómeno en el 14,4\,\% de las piezas,
por encima de todas las anotadoras salvo una (UCM3-4, con el 27,1\,\%). También
\texttt{gemini} marca \texttt{masc\_generico} (66,9\,\%) dentro del rango humano. La ceguera global de los modelos convive, por tanto, con
puntos de mayor sensibilidad en variables concretas, habitualmente las más explícitas y sencillas de definir.

Hay además un segundo dato en la misma tabla: la dispersión \emph{entre las propias personas} es mayor que la que
separa a los modelos entre sí. La prevalencia media humana va de 0,283 a 0,639 (un
rango de 0,356), mientras que la de los cuatro modelos va de 0,043 a 0,183 (un rango
de 0,14). El criterio humano no es un punto de referencia fijo, sino una franja
ancha, y esa amplitud es justamente la que acota el acuerdo alcanzable. Como cada persona
codificó piezas distintas, su prevalencia mezcla dos factores: el criterio propio y
la muestra que le tocó.

Un segundo análisis mide el acuerdo \emph{entre} modelos, que sí tienen solapamiento
total porque los cuatro codificaron las 1.313 piezas. La
Tabla~\ref{tab:iris-intraia} lo recoge variable a variable, sin promediar, ya que
los coeficientes de variables con prevalencias tan distintas no son comparables
entre sí (Sección~\ref{subsec:iris_divergencia}).

% Generado por gen_tabla_intraia.py. No editar a mano.
\begin{table}[H]
    \centering
    \begingroup
    \scriptsize
    \setlength{\tabcolsep}{4pt}
    \renewcommand{\texttt}[1]{{\ttfamily #1}}
    \ttabbox[\FBwidth]{
        \caption{Acuerdo entre modelos ($\kappa$ de Cohen) por variable y nivel ($N=1.313$). Cada celda muestra B0 / B1. No se promedia entre variables, porque sus prevalencias no lo permiten.}
        \label{tab:iris-intraia}
    }{
        \begin{tabularx}{\textwidth}{
            >{\hsize=1.6\hsize\raggedright\arraybackslash}X
            >{\hsize=0.88\hsize\centering\arraybackslash}X
            >{\hsize=0.88\hsize\centering\arraybackslash}X
            >{\hsize=0.88\hsize\centering\arraybackslash}X
            >{\hsize=0.88\hsize\centering\arraybackslash}X
            >{\hsize=0.88\hsize\centering\arraybackslash}X
        }
            \toprule
            \textbf{Par de modelos} & \textbf{V25} & \textbf{V26} & \textbf{V30} & \textbf{V33} & \textbf{V35} \\
            \midrule
            \texttt{gemini} $\leftrightarrow$ \texttt{gpt-4o-mini} & 0,063 / 0,047 & 0,139 / 0,055 & 0,091 / 0,015 & 0,010 / 0,086 & 0,092 / 0,062 \\
            \texttt{gemini} $\leftrightarrow$ \texttt{gpt-5.4-nano} & 0,056 / $-$0,001 & 0,143 / 0,012 & 0,079 / 0,037 & 0,014 / 0,022 & 0,042 / $-$0,001 \\
            \texttt{gpt-4o-mini} $\leftrightarrow$ \texttt{gpt-5.4-nano} & 0,177 / 0,002 & 0,163 / 0,071 & 0,550 / 0,300 & 0,189 / 0,048 & 0,316 / $-$0,001 \\
            \midrule
            \texttt{gemma} $\leftrightarrow$ \texttt{gemini} & 0,025 / 0,042 & 0,260 / 0,190 & 0,265 / 0,236 & 0,063 / 0,259 & 0,031 / 0,054 \\
            \texttt{gemma} $\leftrightarrow$ \texttt{gpt-4o-mini} & 0,065 / 0,012 & 0,121 / 0,140 & 0,094 / 0,107 & 0,154 / 0,102 & 0,242 / 0,160 \\
            \texttt{gemma} $\leftrightarrow$ \texttt{gpt-5.4-nano} & 0,077 / $-$0,001 & 0,119 / 0,047 & 0,125 / 0,055 & 0,125 / 0,028 & 0,160 / $-$0,001 \\
            \bottomrule
        \end{tabularx}
    }
    \endgroup
\end{table}


El acuerdo entre modelos resulta bajo en casi todas las celdas, aunque con
diferencias apreciables entre variables que el promedio habría ocultado. En el
nivel B1 la mediana de los treinta valores se sitúa en 0,048 y el rango va de
$-0,001$ a 0,300. Los dos únicos casos con algo de convergencia son los dos modelos
de OpenAI en \texttt{sexismo\_discurso} ($\kappa = 0,300$), previsible por proceder
de la misma empresa proveedora, y el par formado por \texttt{gemma4:e4b} y
\texttt{gemini-3.1-flash-lite} en \texttt{asimetria\_mujer\_hombre} (0,259) y en
\texttt{sexismo\_discurso} (0,236), ambos de la familia de modelos de Google. Fuera
de esas coincidencias por familia, los modelos no comparten criterio.

Conviene señalar que el acuerdo bruto entre modelos sí es alto en V25, V33 y V35,
por encima del 92\,\% en la mayoría de los pares, pero se trata de un artefacto de
que casi todos responden \textit{No} en casi todas las piezas. Es la misma cautela
que exige la exactitud frente al clasificador trivial
(Sección~\ref{sec:iris_variables}), aplicada ahora al acuerdo entre sistemas.

No existe, por tanto, un criterio común a los modelos. Cada uno presenta su propio
patrón de respuesta, y lo que comparten se limita a la dirección del sesgo, la
tendencia a marcar de menos, no a los casos concretos en que lo aplican. Este
resultado matiza la idea de un sesgo único y uniforme de la inteligencia artificial,
ya que se observan más bien cuatro comportamientos distintos que solo coinciden en
sus limitaciones de detección.

El patrón se mantiene al repetir el análisis sobre el nivel de control (B0), lo que
indica que procede del modelo y no del empaquetado como \textit{skills}. La
arquitectura introduce, de hecho, un matiz en sentido contrario al esperable, ya que
las \textit{Agent Skills} \textbf{reducen} el acuerdo entre modelos en la mayoría de
las celdas. El descenso más llamativo es el de los dos modelos de OpenAI en
\texttt{sexismo\_discurso}, de 0,550 a 0,300, y en \texttt{denominacion\_sexualizada},
de 0,316 a $-0,001$. Las \textit{skills} acentúan las particularidades de cada
modelo en lugar de acercarlos a un criterio compartido.

Los tres análisis permiten repartir el desajuste entre sus dos causas. El primero
mostró que el acuerdo depende del equipo de anotación tanto como del modelo, lo que
sitúa una parte del desacuerdo en la propia tarea, cuyo criterio no está
suficientemente estabilizado. El segundo mostró que, con independencia de eso, los
modelos comparten una infradetección sistemática. El tercero, que ni siquiera
convergen entre sí, de modo que esa infradetección no responde a un criterio
alternativo compartido.

La lectura, por tanto, tiene dos mitades y ninguna anula a la otra. Existe un
\textbf{techo de la tarea}, visible en la distancia entre los dos equipos de
anotación, que ningún sistema podría superar mientras el criterio admita
interpretaciones tan distintas, y que aconseja no leer un $\kappa$ bajo como medida
directa de incompetencia. Pero existe también una \textbf{limitación real de los
modelos}, que en variables como \texttt{lenguaje\_sexista} no detectan un fenómeno
que ambos equipos coinciden en marcar en la mayoría de las piezas. El peso relativo
de cada mitad cambia según la variable, y su distinción es la que ordena las
implicaciones para el proyecto que se recogen en la
Sección~\ref{sec:discussion}.
